\PassOptionsToPackage{table,xcdraw}{xcolor}
\documentclass[11pt, a4paper]{article}

\usepackage[T1]{fontenc}
\usepackage[utf8]{inputenc}
\usepackage{tgpagella}
\usepackage[spanish, es-noshorthands]{babel}
\usepackage{amsmath, amssymb}
\usepackage[most]{tcolorbox}
\usepackage{tabularx}
\usepackage[margin=2.5cm]{geometry}
\usepackage{fancyhdr}

\pagestyle{fancy}
\fancyhf{}
\setlength{\headheight}{16pt}
\lhead{\textbf{UCB Sede Tarija} | Ing. Mecatrónica}
\rhead{IMT-221 | Hoja de Medición DC S07-3}
\renewcommand{\headrulewidth}{0.4pt}
\definecolor{ucbblue}{RGB}{0, 51, 102}

\begin{document}

\begin{tcolorbox}[colback=ucbblue!5, colframe=ucbblue, title=\textbf{Registro de Mediciones Estáticas (DC) — Par Diferencial}]
    Complete esta tabla con sus predicciones analíticas, la simulación de LTspice y las mediciones físicas en placa antes de proceder al análisis de señal diferencial[cite: 18].
\end{tcolorbox}

\section*{1. Valores de Componentes y Rieles}
\renewcommand{\arraystretch}{1.6}
% Se utiliza \raggedright y \centering para evitar advertencias de Underfull/Overfull \hbox
\begin{tabularx}{\textwidth}{|>{\raggedright\arraybackslash}X|>{\centering\arraybackslash}X|>{\centering\arraybackslash}X|}
\hline
\rowcolor{ucbblue!20}
\textbf{Componente / Riel} & \textbf{Valor Nominal[cite: 18]} & \textbf{Valor Real Medido} \\ \hline
Resistor $R_{C1}$ & $4.7\,\text{k}\Omega$ & \\ \hline
Resistor $R_{C2}$ & $4.7\,\text{k}\Omega$ & \\ \hline
Resistor de Referencia $R_{REF}$ & $8.2\,\text{k}\Omega$ & \\ \hline
Resistor Limitador $R_{LIM}$ & $1.8\,\text{k}\Omega$ & \\ \hline
Resistor Shunt $R_{sh}$ & $1\,\Omega$ & \\ \hline
Voltaje Positivo $V_{CC}$ (vs $0\,\text{V}$) & $+9.0\,\text{V}$ & \\ \hline
Voltaje Negativo $V_{EE}$ (vs $0\,\text{V}$) & $-9.0\,\text{V}$ & \\ \hline
\end{tabularx}

\section*{2. Subsistema: Espejo de Corriente (Cola)}
\begin{tabularx}{\textwidth}{|>{\raggedright\arraybackslash}p{4.5cm}|>{\centering\arraybackslash}X|>{\centering\arraybackslash}X|>{\centering\arraybackslash}X|}
\hline
\rowcolor{ucbblue!20}
\textbf{Variable DC} & \textbf{Predicho Analíticamente} & \textbf{Simulado (LTspice)} & \textbf{Medido en Hardware} \\ \hline
$V_{REF}$ (Voltaje Base Q3/Q4) & $\approx -8.3\,\text{V}$[cite: 18] & & \\ \hline
$I_{REF}$ (Corriente en $R_{REF}$) & $\approx 1.01\,\text{mA}$[cite: 18] & & \\ \hline
$I_T$ (Corriente de Cola) & $\approx 1\,\text{mA}$[cite: 18] & & \\ \hline
\end{tabularx}

\vspace{0.3cm}
\textit{Nota: No conecte Q1 ni Q2 hasta validar que la corriente de cola $I_T$ está en el orden de $1\,\text{mA}$[cite: 18].}

\section*{3. Subsistema: Par Diferencial (Estado Balanceado, $v_1=v_2=0\,\text{V}$)}
\begin{tabularx}{\textwidth}{|>{\raggedright\arraybackslash}p{4.5cm}|>{\centering\arraybackslash}X|>{\centering\arraybackslash}X|>{\centering\arraybackslash}X|}
\hline
\rowcolor{ucbblue!20}
\textbf{Variable DC} & \textbf{Predicho Analíticamente[cite: 18]} & \textbf{Simulado (LTspice)} & \textbf{Medido en Hardware} \\ \hline
Voltaje de Emisor ($v_E$) & $\approx -0.7\,\text{V}$ & & \\ \hline
Voltaje Colector 1 ($v_{C1}$) & $\approx 6.65\,\text{V}$ & & \\ \hline
Voltaje Colector 2 ($v_{C2}$) & $\approx 6.65\,\text{V}$ & & \\ \hline
Corriente Colector 1 ($I_{C1}$) & $\approx 0.5\,\text{mA}$ & & \\ \hline
Corriente Colector 2 ($I_{C2}$) & $\approx 0.5\,\text{mA}$ & & \\ \hline
Voltaje Unión $V_{CE1}$ & $\approx 7.35\,\text{V}$ & & \\ \hline
Voltaje Unión $V_{CE2}$ & $\approx 7.35\,\text{V}$ & & \\ \hline
Voltaje de Offset ($V_{OD,0}$) & $0\,\text{V}$ (Ideal) & & \\ \hline
\end{tabularx}

\vspace{0.6cm}
\textbf{Notas de Troubleshooting (Desviaciones observadas, Offset térmico, problemas de Mismatch y correcciones aplicadas)[cite: 18]:}
\vspace{0.3cm}
\hrule\vspace{0.7cm}
\hrule\vspace{0.7cm}
\hrule

\end{document}
